\chapter{Gli infortuni nel calcio}

L'argomento si apre con la lettura dei dati epidemiologici che la ricerca internazionale mette a
disposizione ogni anno: il \textbf{UEFA Elite Club Injury Study Report} nella sua edizione
\textbf{2016/17}, l'\textit{End of Season Injury Review} della Premier League della stessa stagione
e alcuni lavori sui fattori di rischio e sulla prevenzione. La prima metà del capitolo descrive
\textbf{che cosa succede} a livello di infortunio; la seconda affronta il \textit{che cosa fare} per
allontanarne l'incidenza --- i principi della prevenzione, il percorso di rientro, i programmi di
allenamento preventivo, lo screening funzionale e il monitoraggio individuale del giocatore.

% ---------------------------------------------------------------------------
\section{L'UEFA Elite Club Injury Study Report 2016/17}

\subsection{Il campione osservato}

\begin{itemize}
  \item \textbf{21 Club} partecipanti alla UEFA Champions League;
  \item \textbf{170.000 h} di lavoro analizzate, così ripartite:
    \begin{itemize}
      \item 145.000 h di \textbf{allenamento} ($85\%$);
      \item 25.000 h di \textbf{partita} ($15\%$);
    \end{itemize}
  \item per singolo Club, indicativamente \textbf{7500 h} di allenamento ($88\%$) e \textbf{1000 h}
    di partite ($12\%$);
  \item in media ogni giocatore ha sostenuto \textbf{232 sessioni allenanti} e \textbf{60 partite}
    nell'intero periodo, cioè \textbf{21 sessioni} e \textbf{5,5 partite} al mese.
\end{itemize}

\subsection{Il conteggio degli infortuni}

\rubrica{Totali}
\begin{itemize}
  \item \textbf{795 infortuni} complessivi, di cui:
    \begin{itemize}
      \item \textbf{339 in allenamento} ($43\%$);
      \item \textbf{456 in partita} ($57\%$);
    \end{itemize}
  \item nello specifico: \textbf{142 gravi} ($18\%$), \textbf{359 muscolari} ($45\%$),
    \textbf{132 ai legamenti} ($17\%$).
\end{itemize}

\rubrica{Localizzazione}
La distribuzione per distretto separa nettamente l'allenamento dalla gara: in allenamento domina il
comparto \textbf{polpaccio--tendine d'Achille}, in gara il quadro è molto più uniforme.

\begin{table}[htbp]
  \centering
  \small
  \renewcommand{\arraystretch}{1.3}
  \begin{tabularx}{\textwidth}{|X|c|c|}
    \hline
    \textbf{Localizzazione} & \textbf{Allenamento} & \textbf{Gara} \\
    \hline
    Anca--adduttori & 20\% & 13\% \\
    \hline
    Coscia & 27\% & 29\% \\
    \hline
    Ginocchio & 14\% & 17\% \\
    \hline
    Polpaccio--tendine d'Achille & 44\% & 8\% \\
    \hline
    Caviglia & 10\% & 15\% \\
    \hline
  \end{tabularx}
\end{table}

\rubrica{Tipologia}
\begin{itemize}
  \item \textbf{distrazioni e rotture muscolari}: $51\%$ in allenamento, $40\%$ in gara;
  \item \textbf{distorsioni legamentose}: $10\%$ in allenamento, $20\%$ in gara.
\end{itemize}

\subsection{Il meccanismo dell'infortunio}

\begin{table}[htbp]
  \centering
  \small
  \renewcommand{\arraystretch}{1.3}
  \begin{tabularx}{\textwidth}{|X|c|c|}
    \hline
    \textbf{Meccanismo} & \textbf{Allenamento} & \textbf{Gara} \\
    \hline
    Corsa--sprint & 19,4\% & 19,7\% \\
    \hline
    Cambio di direzione (CoD) & 7,4\% & 7,9\% \\
    \hline
    Calciando & 12,7\% & 3,4\% \\
    \hline
    \textit{Tackle} & 6,3\% & 16,2\% \\
    \hline
    \textit{Overuse} & 23,6\% & 9,3\% \\
    \hline
  \end{tabularx}
\end{table}

\begin{itemize}
  \item \textbf{corsa--sprint} e \textbf{cambi di direzione}: valori praticamente sovrapponibili nei
    due contesti;
  \item \textbf{calciando}: nettamente più frequente in allenamento ($12,7\%$ contro $3,4\%$);
  \item \textbf{tackle}: prevalentemente in gara ($16,2\%$ contro $6,3\%$), com'è atteso da un gesto
    di contrasto;
  \item \textbf{overuse}: soprattutto in allenamento ($23,6\%$ contro $9,3\%$).
\end{itemize}

\subsection{Gli infortuni da \textit{overuse}}

Si verificano sulle seguenti strutture:

\begin{itemize}
  \item \textbf{tendini}: tendinopatie;
  \item \textbf{strutture muscolo-tendinee}, in particolare il \textbf{pube};
  \item \textbf{tessuti molli}: sindrome compartimentale da sforzo cronico, caratterizzata da
    crampi, tensione e dolore;
  \item \textbf{sindrome da stress tibiale mediale} (periostite);
  \item \textbf{borsite}: infiammazione da sovraccarico, p.\,es.\ alla caviglia;
  \item \textbf{fascite plantare}, a carico dell'arco plantare;
  \item \textbf{fratture da stress}.
\end{itemize}

\subsection{La gravità e la ricaduta}

\begin{table}[htbp]
  \centering
  \small
  \renewcommand{\arraystretch}{1.3}
  \begin{tabularx}{\textwidth}{|X|c|c|}
    \hline
    \textbf{Giorni di assenza} & \textbf{Allenamento} & \textbf{Gara} \\
    \hline
    1--3 giorni & 16,4\% & 9,4\% \\
    \hline
    4--7 giorni & 27,4\% & 21,4\% \\
    \hline
    8--28 giorni & 41\% & 49\% \\
    \hline
    $>$ 28 giorni & 14\% & 20\% \\
    \hline
  \end{tabularx}
\end{table}

\begin{itemize}
  \item la slide evidenzia due valori aggregati sulle righe centrali: \textbf{48,4\%} per la fascia
    \textbf{4--7 giorni} e \textbf{90\%} per la fascia \textbf{8--28 giorni};
  \item la fascia \textbf{8--28 giorni} è quella che raccoglie gli \textbf{infortuni gravi} ed è la
    più rappresentata in entrambi i contesti.
\end{itemize}

\rubrica{Ricaduta dall'infortunio}
\begin{itemize}
  \item \textbf{nessuna ricaduta}: $91,3\%$ in allenamento, $92,9\%$ in gara;
  \item \textbf{ricaduta}: $8,4\%$ in allenamento, $6,4\%$ in gara.
\end{itemize}

\subsection{Incidenza, assenze e distribuzione nella stagione}

\rubrica{Allenamento}
\begin{itemize}
  \item incidenza media \textbf{2,3 infortuni ogni 1000 h} di allenamento, con indice di squadra da
    \textbf{0,1} fino a \textbf{4,2};
  \item assenza media \textbf{16 giorni}, con indice di squadra da \textbf{8} a \textbf{31} giorni;
  \item la distribuzione mensile degli infortuni in allenamento si concentra nella parte centrale
    del girone di andata, con i valori più alti in \textbf{settembre--ottobre} --- riconducibili
    all'addensarsi degli impegni fra campionato e coppe --- e un calo in \textbf{novembre} e
    soprattutto in \textbf{dicembre}, per la pausa invernale o la riduzione del numero di partite.
\end{itemize}

\rubrica{Gara}
\begin{itemize}
  \item incidenza media \textbf{19,8 infortuni ogni 1000 h} di gioco, con indice di squadra da
    \textbf{7,1} a \textbf{37,6};
  \item assenza media \textbf{23 giorni}, con indice di squadra da \textbf{7} a \textbf{44} giorni;
  \item la distribuzione mensile degli infortuni in partita --- riferita alla squadra campione
    (\textit{Team X}) --- è più elevata in \textbf{agosto} (per i carichi della preparazione?) e nei
    mesi di \textbf{gennaio} e \textbf{febbraio} (condizioni meteo? maggior numero di partite
    giocate?): sono ipotesi interpretative, non conclusioni dello studio.
\end{itemize}

\rubrica{Ricadute e conseguenze}
\begin{itemize}
  \item in media solo l'\textbf{8\%} degli infortuni subiti ha avuto una ricaduta, con indice di
    squadra da $0\%$ a $17\%$: il dato indica un buon lavoro degli staff medici nella gestione del
    recupero;
  \item fra le \textbf{ragioni di assenza} dagli allenamenti e dalle gare, l'infortunio è di gran
    lunga la voce prevalente; malattia, altri motivi e impegni con le nazionali pesano molto meno;
  \item in media ogni giocatore ha perso \textbf{1,9 allenamenti} e \textbf{0,5 partite} al mese a
    causa degli infortuni.
\end{itemize}

% ---------------------------------------------------------------------------
\section{La Premier League: \textit{End of Season Injury Review} 2016--17}

\subsection{Il costo della stagione}

\begin{itemize}
  \item \textbf{20.576 giorni} complessivamente persi per infortunio;
  \item \textbf{614 infortuni significativi}, cioè con assenza di \textbf{10 giorni o più};
  \item \textbf{£131.314.980} di monte ingaggi pagato a giocatori infortunati;
  \item per singolo Club gli infortuni significativi vanno da \textbf{47} (Sunderland, 1813 giorni
    persi) a \textbf{14} (West Bromwich Albion, 319 giorni persi).
\end{itemize}

\subsection{Gli infortuni per distretto corporeo}

Su un totale di \textbf{1111} infortuni registrati:

\begin{multicols}{2}
\begin{itemize}
  \item \textit{hamstring} (flessori): \textbf{166} --- il $15\%$ di tutti gli infortuni;
  \item caviglia--piede: \textbf{161};
  \item ginocchio: \textbf{151}, di cui \textbf{12} rotture del legamento crociato anteriore;
  \item non specificati: \textbf{149} (108 \textit{knocks}, 41 muscolari);
  \item inguine--anca: \textbf{115};
  \item polpaccio--tibia: \textbf{112};
  \item malattia: \textbf{89};
  \item coscia: \textbf{56};
  \item schiena: \textbf{36};
  \item testa--viso: \textbf{35};
  \item torace--addome: \textbf{14};
  \item spalla: \textbf{14};
  \item polso--mano: \textbf{8};
  \item braccio--gomito: \textbf{3};
  \item collo: \textbf{2}.
\end{itemize}
\end{multicols}

I \textbf{flessori della coscia} sono il distretto più colpito, seguiti dalla caviglia e dal
ginocchio.

\subsection{Quando avvengono e chi colpiscono}

\rubrica{Momento della gara}
\begin{itemize}
  \item riscaldamento: $2\%$;
  \item primo tempo: $24\%$;
  \item intervallo: $8\%$;
  \item secondo tempo: $66\%$.
\end{itemize}

Il picco è al \textbf{90$'$} --- 26 giocatori sostituiti in quel minuto, recupero incluso ---
seguito dal \textbf{57$'$} e dal \textbf{60$'$}, con 16 sostituzioni per infortunio ciascuno.

\rubrica{Ruolo in campo}
\begin{itemize}
  \item portieri: $5\%$;
  \item difensori: $34\%$;
  \item centrocampisti: $39\%$;
  \item attaccanti: $22\%$.
\end{itemize}

Il quadro che ne esce, incrociato con i dati UEFA, è che gli infortuni si concentrano
\textbf{in gara} ($57\%$) e sono in prevalenza \textbf{muscolari} ($45\%$).

% ---------------------------------------------------------------------------
\section{Il carico come predittore dell'infortunio}

Molti Club considerano il monitoraggio della \textbf{High Speed Running} (HSR) un indicatore del
carico. In Premier League:

\begin{itemize}
  \item gli \textbf{sprint} coprono solo il \textbf{10\%} della distanza totale di gara;
  \item \textbf{accelerazioni e decelerazioni} avvengono ogni \textbf{9 secondi} nei 90 minuti;
  \item in media il calciatore effettua \textbf{656 accelerazioni} ($> 0,5$ m/s$^2$) e
    \textbf{612 decelerazioni} ($< -0,5$ m/s$^2$) in una partita;
  \item diverse ricerche attuali indicano che \textbf{accelerazioni e decelerazioni} pesano sulla
    richiesta fisiologica più della \textbf{corsa in linea}.
\end{itemize}

\subsection{Gli scostamenti dallo standard individuale (Kitman Labs)}

Il rischio di infortunio si lega allo \textbf{scostamento del carico degli ultimi 7 giorni} rispetto
alla \textbf{media individuale} dell'atleta, misurato in deviazioni standard (SD).

\begin{table}[htbp]
  \centering
  \small
  \renewcommand{\arraystretch}{1.3}
  \begin{tabularx}{\textwidth}{|X|c|c|}
    \hline
    \textbf{Parametro (ultimi 7 giorni)} & \textbf{Scostamento} & \textbf{Rischio di infortunio} \\
    \hline
    High Speed Running & $-1,5$ SD & $+57\%$ \\
    \hline
    Accelerazioni totali & $+1,25$ SD & $\times 1,8$ \\
    \hline
    Decelerazioni totali & $-1$ SD & $+25\%$ \\
    \hline
    Cambi di direzione & $-1,75$ SD & $+88\%$ \\
    \hline
    Distanza totale & $-1,75$ SD & $+86\%$ \\
    \hline
  \end{tabularx}
\end{table}

Il dato controintuitivo è che quasi tutti questi parametri aumentano il rischio \textbf{quando
diminuiscono}: è il \textbf{sotto-carico} rispetto all'abitudine individuale a esporre il giocatore,
non solo il sovraccarico. Fa eccezione l'accelerazione totale, che pesa in \textbf{aumento}.

\subsection{Le implicazioni operative}

\begin{itemize}
  \item il \textbf{decremento del numero di decelerazioni} rispetto allo standard individuale
    dell'atleta incrementa il rischio di infortunio in gara;
  \item esporre il calciatore a un \textbf{eccessivo carico di allenamento mirato alle
    decelerazioni} potrebbe però provocare un decremento nella performance agonistica;
  \item in ordine all'ottimizzazione della performance sull'\textbf{intera stagione}, gli staff
    dovrebbero considerare come accelerazioni e decelerazioni contribuiscano al \textbf{sovraccarico
    fisiologico} e alla \textbf{fatica} dell'atleta.
\end{itemize}

% ---------------------------------------------------------------------------
\section{I fattori di rischio}

\subsection{Fattori legati all'atleta}

\begin{itemize}
  \item \textbf{età del calciatore}: con l'avanzare dell'età l'usura delle strutture
    muscolo-tendinee diventa un fattore di peso;
  \item \textbf{carriera} (anni di attività agonistica): conta anche il \textit{tipo} di carriera
    --- il solo campionato nazionale espone molto meno della somma di campionato, coppe europee e
    attività con le nazionali;
  \item \textbf{instabilità funzionale}: flessori, anche, ginocchia;
  \item \textbf{instabilità meccanica}: caviglia e ginocchia;
  \item \textbf{recidiva da infortuni}: gli atleti con precedenti infortuni si dimostrano a rischio
    di recidiva \textbf{nell'arco della stessa stagione} --- il recupero incompleto è la condizione
    che espone di più.
\end{itemize}

\subsection{Le superfici di gioco}

\begin{itemize}
  \item secondo alcune ricerche il rischio di infortunio \textbf{non è maggiore} sui terreni
    artificiali;
  \item gli \textbf{infortuni muscolari al comparto coscia} sono \textbf{minori} su terreno
    artificiale;
  \item le \textbf{distorsioni alla caviglia} hanno invece \textbf{maggiore rilevanza}: la
    differenza è riconducibile al \textit{grip} del piede all'impatto, maggiore sul terreno
    artificiale, più rigido, rispetto al terreno erboso, più morbido.
\end{itemize}

\rubrica{Lo studio di Ekstrand, Hägglund e Fuller (2010)}
\textit{Comparison of injuries sustained on artificial turf and grass by male and female elite
football players.}

\begin{itemize}
  \item \textbf{obiettivo}: confrontare incidenze e tipologie di infortunio di squadre di élite
    maschili e femminili su \textbf{erba sintetica} e \textbf{erba naturale};
  \item \textbf{campione}: 20 squadre (15 maschili, 5 femminili);
  \item \textbf{esposizione}: \textbf{246.000 h} complessive --- 143.571 h su sintetico e 54.500 h
    su erba per gli uomini, 38.555 h su sintetico e 9849 h su erba per le donne;
  \item \textbf{2105 infortuni} registrati, di cui il \textbf{71\%} traumatici e il \textbf{29\%} da
    \textit{overuse} (proporzione identica nei due sessi);
  \item \textbf{risultato}: nessuna differenza significativa nelle lesioni da \textit{overuse} fra
    erba sintetica ed erba naturale, né per gli uomini né per le donne.
\end{itemize}

\subsection{Stile di vita e fattori psicofisici}

\begin{itemize}
  \item \textbf{non corretto stile di vita};
  \item \textbf{disturbi del sonno}: gli atleti che dormono in media \textbf{meno di 8 ore} per
    notte hanno un rischio di infortunio \textbf{1,7 volte maggiore} rispetto a chi ne dorme
    \textbf{8 o più}; la probabilità di infortunio su 21 mesi cresce al ridursi delle ore di sonno e
    cala nettamente a partire dalle 8 ore;
  \item \textbf{aspetti nutrizionali}: alimentazione corretta e idratazione;
  \item \textbf{percezione della fatica};
  \item \textbf{eccessivo minutaggio di gioco}, che comporta un \textit{training load} individuale
    elevato;
  \item \textbf{stress psicologici}, da non mettere mai in secondo piano.
\end{itemize}

% ---------------------------------------------------------------------------
\section{I principi della prevenzione}

\subsection{Che cosa fare per prevenire}

\begin{itemize}
  \item \textbf{condizionamento fisico adeguato}: uno stato di salute idoneo alla partecipazione
    all'attività calcistica;
  \item \textbf{non trascurare riscaldamento e defaticamento};
  \item \textbf{controllo dei carichi di lavoro}, a volte non corretti o eccessivi;
  \item \textbf{gesto tecnico corretto}: una mancata coordinazione nell'esecuzione tecnica può
    tradursi in problematiche muscolari, legamentose o tendinee;
  \item \textbf{abbigliamento sportivo adeguato}, e in concreto l'uso delle attrezzature
    fondamentali, i parastinchi in primo luogo;
  \item \textbf{screening medico preventivo appropriato}: conoscere lo \textbf{storico
    dell'atleta} --- prestazioni, prerequisiti fisici e soprattutto le problematiche ricorrenti
    negli anni --- per tenere sotto controllo il progetto individuale;
  \item \textbf{infortuni precedenti non adeguatamente trattati}: a fine stagione, dopo 10--11 mesi
    di attività agonistica e con un indice di \textit{overuse} accumulato, manca quasi sempre un
    percorso che riequilibri ciò che la stagione ha squilibrato --- gli infortuni occorsi, ma anche
    le carenze di flessibilità ed elasticità muscolare emerse;
  \item \textbf{mancato rispetto delle regole del gioco};
  \item \textbf{poco tempo dedicato agli esercizi di prevenzione}: le giustificazioni abituali ---
    mancanza di locali adatti, l'allenatore che vuole subito la squadra in campo --- non reggono.
\end{itemize}

\subsection{Una prevenzione multifattoriale}

Deve essere \textbf{ad ampio raggio} e interessare più elementi:

\begin{itemize}
  \item \textbf{preparazione fisica}, in tutte le aree che la interessano;
  \item \textbf{equipaggiamento}:
    \begin{itemize}
      \item \textbf{parastinchi}: traumi diretti su tibia e perone;
      \item \textbf{bendaggio delle caviglie}: re-infortunio da distorsione della caviglia;
    \end{itemize}
  \item \textbf{fair play} fra gli atleti, sia in allenamento sia in gara;
  \item \textbf{percorsi diagnostico-terapeutico-riabilitativi}: routine di controllo, farmaci,
    chirurgia, fisioterapia;
  \item \textbf{riatletizzazione accurata}, dove entra in gioco il preparatore fisico in
    relazione con lo staff medico e fisioterapico.
\end{itemize}

\subsection{Mezzi attivi e mezzi passivi}

\begin{itemize}
  \item \textbf{mezzi attivi}: iniziative che l'\textbf{atleta} mette in atto per proteggersi
    efficacemente dall'azione potenzialmente dannosa dell'esercizio fisico:
    \begin{itemize}
      \item corretta esecuzione dei movimenti;
      \item esatta percezione del proprio corpo, cioè sentire il corpo che si muove nelle varie
        direzioni e situazioni;
      \item adeguato grado di allenamento;
    \end{itemize}
  \item \textbf{mezzi passivi}: precauzioni adottate direttamente \textbf{sull'atleta e
    sull'ambiente}, con lo scopo di eliminare o quanto meno ridurre il rischio nelle varie attività
    sportive.
\end{itemize}

% ---------------------------------------------------------------------------
\section{La prevenzione: il programma FIFA 11+}

\subsection{Che cos'è l'11+}

\begin{itemize}
  \item nel \textbf{2003} l'\textbf{F-MARC} sviluppò \textbf{``The 11''}, programma di prevenzione
    degli infortuni per calciatori \textbf{dilettanti};
  \item l'efficacia è stata sancita da uno studio condotto in \textbf{Svizzera}:
    l'implementazione su base nazionale (\textbf{2004--2008}) ha determinato un significativo
    decremento degli infortuni \textbf{sia in allenamento sia in partita}, confermando insieme
    all'efficacia anche la \textbf{facilità di impiego} del programma;
  \item dal \textbf{2006} ``The 11'' è stato sviluppato in un programma più completo:
    \textbf{``11+''}.
\end{itemize}

\subsection{La struttura dell'11+}

Tre parti, \textbf{15 esercizi} in totale, da effettuare nella \textbf{sequenza specifica}
all'inizio di ogni seduta di allenamento. Determinanti la \textbf{tecnica corretta} di ciascun
esercizio, la \textbf{postura} e il \textbf{controllo del corpo}.

\begin{enumerate}
  \item \textbf{Parte 1}: esercizi di \textbf{corsa a bassa intensità}, combinati con esercitazioni
    di \textbf{stretching attivo} e di \textbf{contatti controllati} con i compagni;
  \item \textbf{Parte 2}: \textbf{sei serie} di esercizi, ciascuna con \textbf{tre livelli
    progressivi} di difficoltà, mirate a sviluppare la forza dei \textbf{muscoli stabilizzatori
    della zona addominale-pelvica}, l'\textbf{equilibrio}, l'\textbf{abilità pliometrica} e
    l'\textbf{agilità};
  \item \textbf{Parte 3}: esercizi di \textbf{corsa a intensità moderata/alta}, combinati con
    movimenti di \textbf{arresto} e \textbf{cambio di direzione}.
\end{enumerate}

\subsection{L'efficacia: la meta-analisi di Thorborg e coll.\ (2017)}

\textit{Effect of specific exercise-based football injury prevention programmes on the overall
injury rate in football: a systematic review and meta-analysis of the FIFA 11 and 11+ programmes}
(\textit{British Journal of Sports Medicine}).

\rubrica{Impianto}
\begin{itemize}
  \item \textbf{6 studi randomizzati} sull'effetto dei programmi di prevenzione nel calcio
    \textbf{non professionistico};
  \item ripartiti in \textbf{FIFA 11} (2 studi) e \textbf{FIFA 11+} (4 studi).
\end{itemize}

\rubrica{Risultati}
\begin{itemize}
  \item l'\textbf{analisi primaria} mostra una riduzione complessiva degli infortuni, con un
    rapporto di rischio di \textbf{0,75} ($95\%$ CI 0,57--0,98; $p = 0{,}04$) a favore dei programmi
    di prevenzione a protocollo FIFA;
  \item l'\textbf{analisi secondaria}, sui 4 studi che applicano il solo \textbf{FIFA 11+}, rivela
    un'ulteriore riduzione del rapporto di rischio: \textbf{IRR 0,61} ($95\%$ CI 0,48--0,77;
    $p < 0{,}001$).
\end{itemize}

\rubrica{Conclusioni}
\begin{itemize}
  \item il programma \textbf{FIFA 11+} evidenzia una riduzione sostanziale degli infortuni, di circa
    il \textbf{39\%};
  \item riduce le \textbf{prime quattro lesioni più diffuse} nel calcio:
    \begin{itemize}
      \item \textbf{flessori}: $-60\%$;
      \item \textbf{ginocchio}: $-48\%$;
      \item \textbf{adduttori/anche}: $-41\%$;
      \item \textbf{caviglia}: $-32\%$.
    \end{itemize}
\end{itemize}

% ---------------------------------------------------------------------------
\section{Il percorso di rientro dopo l'infortunio}

\subsection{Le cinque fasi}

Dall'infortunio al rientro in squadra il percorso attraversa cinque fasi, ciascuna con una figura
professionale responsabile. Le \textbf{fasi 1 e 2} sono di pertinenza esclusiva dello \textbf{staff
medico}; nelle successive entrano riabilitatori e preparatori fisici.

\begin{enumerate}
  \item \textbf{Fase 1} --- \textit{medico}: primo inquadramento diagnostico e progettazione del
    piano di intervento;
  \item \textbf{Fase 2} --- \textit{fisioterapisti e osteopati}: trattamento medico e fisioterapia;
  \item \textbf{Fase 3} --- \textit{preparatori fisici} per il miglioramento della condizione
    fisica, \textit{riadattatori} per il recupero della struttura lesa e il riadattamento del gesto
    sportivo; si chiude con il \textbf{nulla osta medico} (\textit{medical all-clear});
  \item \textbf{Fase 4} --- \textit{preparatori fisici}: lavoro individuale specifico
    supplementare; si chiude con il \textbf{nulla osta a giocare} (\textit{all-clear to play});
  \item \textbf{Fase 5} --- \textit{allenatori}: rientro pieno in gruppo.
\end{enumerate}

\subsection{Che cosa serve allo staff medico (fasi 1--2)}

\begin{itemize}
  \item \textbf{importanza della clinica}: anamnesi ed esame obiettivo;
  \item formulare un \textbf{sospetto diagnostico preciso};
  \item conoscere \textbf{tutti gli aspetti di una lesione muscolare} e la loro incidenza
    sport-specifica;
  \item conoscere gli \textbf{atti terapeutici} e i \textbf{principi riabilitativi}.
\end{itemize}

\rubrica{L'ordine di grandezza del problema}
\begin{itemize}
  \item in una squadra di 25 giocatori sono previste mediamente, a stagione (Ekstrand, 2011):
    \textbf{15 lesioni muscolari}, \textbf{37 partite perse}, \textbf{223 giorni} di assenza media e
    \textbf{148 giorni} di assenza dagli allenamenti;
  \item ogni campionato i calciatori di un Club perdono, per i soli infortuni ai flessori,
    \textbf{90 allenamenti} e \textbf{15 gare} (Mueller-Wohlfahrt, 2013).
\end{itemize}

\subsection{La recidiva}

\rubrica{Quanto pesa}
\begin{itemize}
  \item il \textbf{30\%} delle lesioni agli \textit{hamstring} ha una recidiva entro
    \textbf{2 settimane} (Heiderscheit, 2010);
  \item il \textbf{59\%} recidiva entro il \textbf{primo mese} (Reurink, 2014);
  \item il tasso di recidiva nel \textbf{primo anno} va dal \textbf{14\%} al \textbf{63\%}
    (Reurink, 2014);
  \item la recidiva è di solito \textbf{più grave} e richiede tempi di recupero più lunghi
    (Reurink, 2014).
\end{itemize}

Le recidive a due settimane e a un mese indicano che staff medico e staff tecnico hanno
\textbf{anticipato i tempi fisiologici} del recupero completo.

\rubrica{Perché avviene}
\begin{itemize}
  \item \textbf{diagnosi sottostimata};
  \item \textbf{programma riabilitativo inadeguato};
  \item \textbf{ritorno prematuro} all'attività sportiva;
  \item \textbf{formazione di tessuto fibrotico}, che sottrae al muscolo l'elasticità;
  \item \textbf{dolore} persistente;
  \item \textbf{inibizione muscolare}, con movimenti non controllati.
\end{itemize}

È di fondamentale importanza \textbf{ottimizzare la diagnosi e il processo terapeutico} per ridurre
il rischio di recidiva.

\subsection{I gradi della lesione muscolare}

\begin{itemize}
  \item \textbf{grado I} (\textit{overstretch}): elongazione, con piccole lacerazioni;
  \item \textbf{grado II}: rottura parziale, più estesa ma incompleta;
  \item \textbf{grado III}: rottura completa del ventre muscolare.
\end{itemize}

\subsection{I tempi di recupero}

Il tempo è \textbf{estremamente variabile e individuale}: la fisiologia è la stessa per tutti, ma la
soglia di sopportazione del dolore cambia da giocatore a giocatore.

\begin{table}[htbp]
  \centering
  \small
  \renewcommand{\arraystretch}{1.3}
  \begin{tabularx}{\textwidth}{|l|l|X|}
    \hline
    \textbf{Fase} & \textbf{Durata} & \textbf{Che cosa avviene} \\
    \hline
    Emostasi & Da 1 a 3 giorni & Fermare l'emorragia \\
    \hline
    Infiammazione & Da 3 a 20 giorni & Nuovo sito di crescita di vasi sanguigni \\
    \hline
    Ricostituzione cellulare & Da 1 a 6 settimane & Chiusura della ferita \\
    \hline
    Ricostruzione del tessuto & Da 6 settimane a 2 anni & Tessuto di riparazione corretto e
      appropriato \\
    \hline
  \end{tabularx}
\end{table}

Elementi chiave del programma riabilitativo sono il \textbf{carico ottimale} e il \textbf{rispetto
degli adattamenti fisiologici}.

% ---------------------------------------------------------------------------
\section{I programmi di allenamento per la prevenzione}

\subsection{Le aree di lavoro}

\begin{itemize}
  \item \textbf{mobilità articolare} e \textbf{flessibilità};
  \item \textbf{propriocezione};
  \item \textbf{equilibrio};
  \item \textbf{stabilizzazione del core}: miglioramento del controllo neuromuscolare attraverso
    esercitazioni di natura statica;
  \item \textbf{potenziamento del core}: attivazione e potenziamento muscolare con attività
    moderatamente dinamiche;
  \item \textbf{efficacia del core}: allenamento attraverso movimenti complessi, mirato allo
    sviluppo delle qualità fisiche;
  \item \textbf{allenamento eccentrico}.
\end{itemize}

\subsection{L'allenamento isoinerziale}

\rubrica{Il principio del volano}
\begin{itemize}
  \item la \textbf{resistenza dell'attrezzo è proporzionale alla forza sviluppata}: non c'è un
    carico prefissato;
  \item il volano viene \textbf{accelerato nella fase concentrica} e \textbf{decelerato nella fase
    eccentrica}, il che rende possibile un \textbf{sovraccarico eccentrico} (Berg e coll., 1994;
    Pozzo, 2008, modificato).
\end{itemize}

\rubrica{Le evidenze}
\begin{itemize}
  \item allenamento in condizioni di \textbf{microgravità} (Rittweger e coll., 2007), nato per il
    recupero degli astronauti al rientro sulla Terra;
  \item \textbf{prevenzione degli infortuni muscolari} (Askling e coll., 2003);
  \item \textbf{ipertrofia muscolare precoce} (Seynnes e coll., 2007);
  \item \textbf{possibilità di modulare l'impegno muscolare} (Tous, 2009).
\end{itemize}

\rubrica{Confronto fra carico inerziale e pesi liberi}
\begin{itemize}
  \item \textbf{forza massima eccentrica}: superiore del \textbf{60\%} con il volano rispetto ai
    pesi liberi;
  \item \textbf{forza massima concentrica}: superiore del \textbf{70\%}.
\end{itemize}

\rubrica{Il limite: il tempo di latenza}
Resta una problematica \textbf{non ancora ben definita}: il \textbf{tempo di latenza}, cioè
l'intervallo in cui il volano deve \textbf{riavvolgersi} prima di poter far ripartire l'attività
meccanica dell'atleta che sta eseguendo l'esercizio.

\rubrica{Diffusione}
Le squadre che usano attrezzatura nHANCE\textsuperscript{TM} comprendono, in \textbf{Primera
División}, Barcelona, Real Madrid, Atlético Madrid, Sevilla, Valencia, Real Sociedad, Zaragoza,
Getafe, Levante e Real Betis; in \textbf{Serie A}, Juventus, Inter, Fiorentina, Lazio, Genoa e
Sampdoria; in \textbf{Premier League}, Arsenal, Liverpool, Manchester United, Tottenham, West Ham,
Bolton e Wolverhampton.

\subsection{La prevenzione nella seduta: due esempi di club}

\rubrica{Juventus --- organizzazione della seduta pre-allenamento}
Una scheda giornaliera elenca i giocatori uno per uno con l'orario di convocazione per ciascuna
attività da svolgere \textbf{prima} della seduta in campo: lavori preventivi, programmi individuali
di attivazione, forza pre-allenamento (con carico ridotto all'avvicinarsi della gara), velocità dei
piedi con le scalette e potenza con balzi sui cubi. L'impostazione è \textbf{individuale}: ogni
atleta ha il proprio programma e il proprio orario.

\rubrica{FC Barcelona --- il \textit{tirante musculador}}
Cinturone elastico per il lavoro di \textbf{quadricipiti, ischiocrurali, lombari e glutei}, usato sia
in regime isometrico sia dinamico. Alla sua introduzione la stampa spagnola attribuisce la
riduzione delle lesioni muscolari della squadra.

\subsection{Il piano di prevenzione del FC Barcelona}

\rubrica{I quattro livelli}

\begin{table}[htbp]
  \centering
  \small
  \renewcommand{\arraystretch}{1.3}
  \begin{tabularx}{\textwidth}{|l|X|c|X|}
    \hline
    \textbf{Livello} & \textbf{Programma} & \textbf{Frequenza} & \textbf{Categorie} \\
    \hline
    1 --- Prevenzione in palestra & Prevenzione di squadra; prevenzione individuale &
      $1\times$ / $2\times$ sett. & Stretching, forza \\
    \hline
    2 --- Prevenzione giornaliera nel riscaldamento & Allenamento eccentrico; propriocezione;
      agilità e coordinazione; velocità di reazione & $1\times$ sett. & Stretching, forza,
      propriocezione, core stability, agilità \\
    \hline
    3 --- Circuiti mirati alla prevenzione & Circuiti senza palla; circuiti con \textit{Bosu};
      circuiti con esercitazioni tecniche; circuiti con giochi di possesso & 1--2$\times$ sett. &
      Stretching, forza, propriocezione, core stability, agilità, attività globali \\
    \hline
    4 --- Esercitazioni specifiche a circuito & Esercizi di forza specifici per il gioco del calcio &
      1--2$\times$ sett. & Forza, propriocezione, agilità \\
    \hline
  \end{tabularx}
\end{table}

\rubrica{Stretching}
Per ogni esercizio: codice, dose minima $\rightarrow$ dose massima, frequenza, collocazione nella
settimana.

\begin{itemize}
  \item \textit{Hamstring}:
    \begin{itemize}
      \item \textbf{St1} --- iperestensione del ginocchio con anteversione del bacino: in piedi, con
        flessione del tronco e anteversione del bacino, estendere il ginocchio;
        $1 \times 3 \rightarrow 1 \times 5$; ogni giorno, pre-allenamento;
      \item \textbf{St2} --- stretching passivo: seduti a gambe tese, flessione del tronco in avanti;
        $1 \times 15''$--$20'' \rightarrow 2 \times 15''$--$20''$; minimo 3 volte a settimana,
        post-allenamento;
      \item \textbf{St3} --- stretching attivo con \textit{tirante musculador}: in piedi con la
        cintura sopra le ginocchia, si scende in flessione percependo lo stiramento posteriore;
        $3 \times 4$--$8 \rightarrow 3 \times 4$--8 con 3\,kg; 1 volta ogni 7 giorni con una gara a
        settimana, ogni 10 giorni con due; sempre, a metà ciclo o dopo la gara;
      \item \textbf{St4} --- stretching della catena posteriore (SGA): supini, gambe tese in
        abduzione contro il muro; $1 \times 3'$--$4'$; 1 volta ogni 7--10 giorni; sempre, a metà
        ciclo o dopo la gara;
    \end{itemize}
  \item \textit{quadricipite}:
    \begin{itemize}
      \item \textbf{St5} --- stretching passivo: in appoggio su una gamba, flessione forzata del
        ginocchio; $1 \times 15''$--$20'' \rightarrow 2 \times 15''$--$20''$; minimo 3 volte a
        settimana, pre/post-allenamento;
    \end{itemize}
  \item \textit{adduttori}:
    \begin{itemize}
      \item \textbf{St6} --- stretching in tensione attiva: in piedi, flessione del tronco e anche
        in abduzione, mettendo in tensione con i gomiti; $1 \times 4''$; ogni giorno,
        pre-allenamento;
      \item \textbf{St7} --- stretching passivo: in ginocchio, ginocchia divaricate, arretrando il
        peso del corpo; su un ginocchio, separando la gamba opposta;
        $1 \times 15''$--$20'' \rightarrow 2 \times 15''$--$20''$; minimo 3 volte a settimana,
        pre/post-allenamento;
      \item \textbf{St8} --- stretching della catena posteriore (SGA): come St4;
        $1 \times 3'$--$4'$; 1 volta ogni 7 giorni con una gara a settimana, ogni 10 giorni con due;
        sempre, a metà ciclo o dopo la gara;
    \end{itemize}
  \item \textit{polpaccio}:
    \begin{itemize}
      \item \textbf{St9} --- stretching passivo: contro il muro, con un piede arretrato e poi lo
        stesso a ginocchio flesso; $1 \times 15''$--$20'' \rightarrow 2 \times 15''$--$20''$; minimo
        3 volte a settimana, pre/post-allenamento.
    \end{itemize}
\end{itemize}

\rubrica{Allenamento della forza}
\begin{itemize}
  \item \textit{ischiocrurali}:
    \begin{itemize}
      \item \textbf{STh1} --- \textit{blue belt}: lavoro eccentrico con varianti in caduta, con
        rotazione e ad angoli di flessione diversi; peso corporeo o peso corporeo con 3\,kg,
        $3 \times 4 \rightarrow 3 \times 8$; 1 volta ogni 7 giorni con una gara a settimana, ogni
        10 giorni con due; prima o dopo le 48 ore dalla gara;
      \item \textbf{STh2} --- \textit{splits with body flexion}: lavoro sugli ischiocrurali a peso
        corporeo, $4 \rightarrow 8$ ripetizioni; una o due volte a settimana; a giorni alterni, mai
        a ridosso della gara;
      \item \textbf{STh3} --- \textit{angel}: lavoro sugli ischiocrurali a peso corporeo,
        $4 \rightarrow 8$ ripetizioni; una o due volte a settimana; a giorni alterni, mai a ridosso
        della gara;
    \end{itemize}
  \item \textit{quadricipiti}:
    \begin{itemize}
      \item \textbf{STq1} --- \textit{squat multifunction}: lavoro sulla tecnica dello squat, con
        carico stabilito in base ai risultati del test di potenza;
        $3 \times 8 \rightarrow 3 \times 10$; una o due volte a settimana; a giorni alterni, mai a
        ridosso della gara;
      \item \textbf{STq2} --- \textit{musculador belt}: lavoro eccentrico con varianti in caduta,
        con rotazione e ad angoli di flessione diversi; peso corporeo o peso corporeo con 3\,kg,
        $3 \times 4 \rightarrow 3 \times 8$; 1 volta ogni 7 giorni con una gara a settimana, ogni
        10 giorni con due; prima o dopo le 48 ore dalla gara.
    \end{itemize}
\end{itemize}

\rubrica{Multi intervention}
Esercizi a corpo libero volti a \textbf{potenziare gli antagonisti} e a \textbf{inibire gli
adduttori}. Tutti $1 \times 30''$ in forma dinamica, una volta a settimana, a fine settimana e di
norma 3 giorni prima della gara.

\begin{itemize}
  \item \textbf{MI1} --- \textit{superman}: in quadrupedia, estensione di braccia e gambe; migliora
    anche la coordinazione;
  \item \textbf{MI2} --- \textit{hip abduction on side}: sul fianco, abduzione dell'arto inferiore
    a ginocchio esteso;
  \item \textbf{MI3} --- \textit{dog}: in quadrupedia, abduzione dell'arto inferiore a ginocchio
    flesso.
\end{itemize}

\rubrica{Core stability}
Tutti $1 \times 30''$, una volta a settimana, a fine settimana e di norma 3 giorni prima della gara.

\begin{itemize}
  \item \textbf{CS1} --- \textit{frontal plank}: ponte frontale in contrazione isometrica del core;
    varianti con instabilità, azioni associate e riduzione degli appoggi;
  \item \textbf{CS2} --- \textit{side plank}: ponte laterale, stesse varianti;
  \item \textbf{CS3} --- \textit{back bridge}: attivazione del core in posizione supina, con
    sinergie muscolari aggiuntive (attivazione dei glutei nel sollevamento del bacino, contrazione
    degli adduttori); in forma dinamica.
\end{itemize}

\rubrica{Propriocezione}
Lavoro su superficie instabile (\textit{Bosu}).

\begin{itemize}
  \item \textbf{Pr1} --- appoggio monopodalico su Bossu, con varianti fino al gesto sportivo:
    rieducazione propriocettiva delle articolazioni dell'arto inferiore; $1 \times 30''$ alternando
    la gamba ogni $5''$; una volta a settimana;
  \item \textbf{Pr2} --- appoggio del ginocchio in instabilità: seduti sul Bossu, rieducazione
    propriocettiva della regione lombo-pelvica; $1 \times 30''$; una volta ogni 15 giorni;
  \item \textbf{Pr3} --- \textit{skipping} laterale in instabilità: rieducazione propriocettiva
    attraverso il gesto dinamico specifico; $1 \times 30''$; una volta ogni 15 giorni.
\end{itemize}

Tutti a fine settimana, di norma 3 giorni prima della gara.

\rubrica{Agilità (lavoro neuromuscolare)}
\begin{itemize}
  \item \textbf{AG1} --- coordinazione sulla scaletta con passaggio di palla: \textit{skipping}
    dentro la scaletta, con movimenti complementari; $1 \times 30''$; due volte a settimana, a fine
    settimana --- di norma 3 giorni prima e un giorno prima della gara;
  \item \textbf{AG2} --- circuito di agilità: circuiti con compiti di agilità diversi, alla massima
    velocità possibile; $1$ serie; due volte a settimana, nel riscaldamento.
\end{itemize}

% ---------------------------------------------------------------------------
\section{Lo screening funzionale: il Functional Movement Screen}

\subsection{La batteria e il punteggio}

Il \textbf{Functional Movement Screen} (FMS) è composto da \textbf{sette test}:

\begin{multicols}{2}
\begin{enumerate}
  \item squat profondo;
  \item superamento dell'ostacolo;
  \item affondo in linea;
  \item mobilità delle spalle;
  \item alzata attiva della gamba tesa;
  \item stabilità del tronco;
  \item stabilità rotatoria.
\end{enumerate}
\end{multicols}

\rubrica{Come si assegna il punteggio}
\begin{itemize}
  \item il \textbf{valore grezzo} serve a differenziare la parte destra dalla sinistra;
  \item il \textbf{punteggio finale} di un test è dato dal \textbf{valore minore} fra i due lati ---
    destra 2 e sinistra 3 danno 2 come punteggio del test;
  \item il totale della scheda è su \textbf{21 punti}.
\end{itemize}

\subsection{L'FMS predice davvero l'infortunio?}

\rubrica{Moran e coll.\ (2017)}
\textit{Do Functional Movement Screen (FMS) composite scores predict subsequent injury? A systematic
review with meta-analysis} (\textit{British Journal of Sports Medicine}): il livello di evidenza
scientifica dell'associazione fra la forza e i punteggi ottenuti con l'FMS \textbf{non è
sufficiente} a supportare l'utilizzo dell'FMS come strumento di previsione degli infortuni.

\rubrica{Chalmers e coll.}
\textit{Asymmetry during preseason Functional Movement Screen testing is associated with injury
during a junior Australian football season}:

\begin{itemize}
  \item i giocatori di football australiano \textit{junior} che mostravano \textbf{asimmetrie} nei
    movimenti durante i test FMS pre-stagionali ebbero \textbf{più probabilità} di infortunarsi
    nella stagione regolare rispetto ai giocatori senza asimmetria;
  \item il \textbf{punteggio di soglia} comunemente riportato di $\leq 14$ \textbf{non} è risultato
    associato all'infortunio in questi atleti.
\end{itemize}

Ciò che conta, quindi, non è il punteggio composito ma l'\textbf{asimmetria fra i due lati}.

% ---------------------------------------------------------------------------
\section{Il monitoraggio individuale del giocatore}

\subsection{Il questionario informativo dell'atleta}

Scheda anagrafica e anamnestica compilata a inizio stagione e aggiornata nel corso dell'anno, per
tenere sotto controllo l'andamento fisico del giocatore. Raccoglie:

\begin{itemize}
  \item dati anagrafici e \textbf{ruolo};
  \item \textbf{interventi chirurgici} subiti, con indicazione della sede;
  \item \textbf{fratture, lussazioni, strappi muscolari, distorsioni, stiramenti}, con sede e data;
  \item se nell'arco della \textbf{stessa stagione} ha subito più di un infortunio muscolare;
  \item \textbf{pubalgia}, \textbf{lombosciatalgia}, \textbf{condropatia del ginocchio}, con
    l'epoca di insorgenza;
  \item \textbf{malocclusione mandibolare} e uso del \textit{bite} in allenamento e in gara;
  \item uso di \textbf{plantari} in allenamento e in gara.
\end{itemize}

\subsection{La tabella di valutazione funzionale}

Griglia che affianca, per ciascun test e per quattro sessioni datate, la
\textbf{condizione ``fit''} --- il giocatore sano a inizio stagione --- alla \textbf{condizione
``unfit''}, quella durante o dopo un infortunio. Il confronto fra le due colonne è il criterio di
rientro.

\rubrica{Test preventivo-riabilitativi}
\begin{itemize}
  \item \textbf{test isometrico} e \textbf{test dinamico}, entrambi su quadricipite, ischiocrurali,
    polpaccio e adduttori.
\end{itemize}

\rubrica{Test condizionali}
\begin{itemize}
  \item \textbf{aspetti metabolici}: capacità aerobica, potenza aerobica, $VO_2max$, soglia
    anaerobica;
  \item \textbf{aspetti neuromuscolari}: relazione forza/velocità, potenza, forza massima, forza
    esplosiva, velocità, forza resistente;
  \item \textbf{aspetti coordinativi}: equilibrio statico a occhi aperti e a occhi chiusi, mobilità
    articolare;
  \item \textbf{data analysis GPS}: percentuale di distanza in sprint ad alta intensità, in
    accelerazione ad alta intensità e in decelerazione ad alta intensità, potenza metabolica,
    percentuale dell'indice anaerobico.
\end{itemize}

\subsection{L'asimmetria fra i due arti come indicatore}

Il test isometrico restituisce la forza di picco dei due lati e la loro differenza percentuale. Le
soglie di lettura sono:

\begin{itemize}
  \item \textbf{0--10\%}: differenze fisiologiche, non problematiche;
  \item \textbf{10--15\%}: ancora normali, ma da tenere sotto controllo;
  \item \textbf{15--20\%}: si consiglia di intervenire per evitare peggioramenti;
  \item \textbf{oltre il 20\%}: si consiglia di intervenire.
\end{itemize}

\rubrica{Due casi sugli adduttori}
Nel test isometrico degli adduttori due atleti presentavano differenziali del \textbf{32\%} e del
\textbf{23\%}: il primo ha poi avuto infortuni ripetuti agli adduttori; il secondo, sottoposto a un
protocollo di lavoro mirato alla prevenzione degli adduttori, ha risolto completamente la sua
problematica.

\subsection{Un caso di ricostruzione del LCA}

Giocatore infortunato in gara il \textbf{17 marzo 2018}, operato al \textbf{legamento crociato
anteriore} il \textbf{22 marzo} con tecnica \textbf{StGr} (semitendinoso-gracile). Arto infortunato:
il \textbf{sinistro}.

\rubrica{Valutazione isometrica del quadricipite}

\begin{table}[htbp]
  \centering
  \small
  \renewcommand{\arraystretch}{1.3}
  \begin{tabularx}{\textwidth}{|X|c|c|c|}
    \hline
    \textbf{Data} & \textbf{Arto destro} & \textbf{Arto sinistro} & \textbf{Differenza} \\
    \hline
    13/07/2017 --- pre-campionato & 636 N & 618 N & $-2,8\%$ \\
    \hline
    09/05/2018 --- 49 giorni post-intervento & 519 N & 161 N & $68,99\%$ \\
    \hline
    22/05/2018 --- 62 giorni post-intervento & 571 N & 253 N & $55,56\%$ \\
    \hline
  \end{tabularx}
\end{table}

\begin{itemize}
  \item il test pre-campionato mostrava un \textbf{equilibrio perfetto} fra i due arti --- il
    $2,8\%$ rientra nelle differenze fisiologiche --- e valori di forza decisamente più alti degli
    attuali;
  \item l'\textbf{arto infortunato} (sinistro) rispetto al proprio valore di luglio 2017:
    $-73,87\%$ a 49 giorni, $-58,94\%$ a 62 giorni;
  \item l'\textbf{arto sano} (destro) rispetto allo stesso riferimento: $-18,36\%$ a 49 giorni,
    $-10\%$ a 62 giorni --- cala anch'esso, per il decondizionamento generale, ma rientra quasi nella
    norma già a 62 giorni;
  \item la differenza fra i due arti resta elevata a 62 giorni, ma il processo di riatletizzazione
    è ben avviato.
\end{itemize}

\rubrica{Valutazione isometrica dei flessori}
Nel test del 22/05/2018 (62 giorni dall'intervento): arto destro \textbf{178 N}, arto sinistro
\textbf{69 N}, differenza \textbf{61,29\%} --- ancora molto elevata.
